\section{SPI på fyra ledningar}
\label{sec:spilines}

\begin{booktable}{@{}llL@{}}
\thead{Ledning} & \thead{Riktning} & \thead{Betydelse} \\ \midrule
\code{SCK} & master $\rightarrow$ slav & Klocka. Ingenting händer utan flanker på den. \\
\code{MOSI} & master $\rightarrow$ slav & Masterns data. \\
\code{MISO} & slav $\rightarrow$ master & Slavens data. \\
\code{SS} & master $\rightarrow$ slav & Aktiv låg. Ramar transaktionen. \\
\end{booktable}

Den viktigaste egenskapen: \textbf{SPI är full duplex.} Varje klockpuls skiftar samtidigt en bit
ut på \code{MOSI} och en bit in på \code{MISO}. Det finns inget sätt att "bara läsa" --- för att få
en byte in måste man skicka en byte ut.

Därför skickar en läsning fyra dummybytes \code{0x00}, och därför är de fyra \code{MISO}-byten
under en skrivning meningslöst skräp.

Parametrarna: SPI mode 0 (klockan vilar låg, data samplas på stigande flank), MSB först,
\code{SCK} högst 1~MHz.

\section{Transaktionen}
\label{sec:transaction}

\textbf{Varje transaktion är exakt fem bytes}: en kommandobyte följd av fyra databytes, med
\code{SS} låg hela vägen.

\begin{console}
Bit:      7    6    5    4    3    2    1    0
        +----+----+----+----+----+----+----+----+
        | W  | 0  | 0  | 0  |   register index  |
        +----+----+----+----+----+----+----+----+
\end{console}

\code{W = 1} betyder skrivning, \code{W = 0} läsning. Bitarna 6--4 är reserverade och ska vara
\code{0}. Bitarna 3--0 är registerindexet, 0--12, ur registerkartan i \kapref{ch:can}.

\subsection{Skrivning}

Mastern skickar registervärdet i de fyra databyten, \textbf{MSB först}. Slaven verkställer det
hopsatta värdet \textbf{först när den femte byten är hel}.

\subsection{Läsning}

I slutet av kommandobyten \textbf{låser slaven registrets aktuella värde en gång}, in i sin
svarsskiftare. De fyra databyten klockar sedan ut det värdet.

\textbf{Låsningsregeln är en del av kontraktet, inte en implementationsdetalj.} De fyra byten hör
alltid till \emph{ett} sammanhängande prov av registret, taget i ett ögonblick. Utan den regeln
hade en \code{STATUS}-läsning kunnat bli ihopsydd av två tidpunkter --- TX-klar-biten från före en
sändning och RX-giltig-biten från efter. Det hade varit en bugg som uppträder ungefär en gång på
tusen och aldrig går att återskapa.

\subsection{\texorpdfstring{\code{SS}}{SS}-reglerna}

Tre regler, och alla tre har kostat någon en kväll:

\begin{itemize}
  \item \textbf{En transaktion per låg \code{SS}-period.} Mastern måste släppa \code{SS} hög
    mellan transaktioner.
  \item \textbf{Back-to-back fungerar inte, och misslyckas tyst.} En sjätte byte som kommer medan
    \code{SS} fortfarande är låg \emph{kastas} --- den tolkas inte som en ny kommandobyte. Att
    slaven är ledig räcker alltså inte.
  \item \textbf{Avbrott har inga sidoeffekter.} Går \code{SS} hög före femte byten överges
    transaktionen helt: ingenting verkställs, ingen trigger går. Det är vad som gör slaven
    självläkande efter en glitch.
\end{itemize}

Det är därför \code{ByteTransport} i \kapref{ch:byte} har \code{select()} och \code{deselect()}
som egna metoder: ramningen är drivrutinens beslut, inte transportens.

\subsection{\texorpdfstring{\code{MISO}}{MISO} utanför en läsning}

Byten mastern klockar ut \textbf{under kommandobyten} är meningslös. Slaven laddar sin
svarsskiftare i \emph{slutet} av kommandobyten, så det som lämnar under den är rester från förra
gången. \textbf{Kasta den.}

\subsection{Ett komplett exempel}

Att skicka en frame (\code{id 0x123}, \code{dlc 2}, data \code{AA BB}):

\begin{console}
Skriv TX_ID      : 81 00 00 01 23
Skriv TX_DLC     : 82 00 00 00 02
Skriv TX_DATA_LO : 83 AA BB 00 00
Skriv TX_DATA_HI : 84 00 00 00 00
Skriv TX_SEND    : 85 00 00 00 01
Läs   STATUS     : 00 xx xx xx xx   -> MISO ger 00 00 00 00 under sändning,
                                      sedan 00 00 00 01 när porten är fri igen.
\end{console}

Läs av kommandobyten i varje rad: \code{0x81} är skrivbiten plus index 1, alltså \code{TX_ID}.
\code{0x00} på sista raden är en läsning av index 0.

Notera \code{83 AA BB 00 00}: \textbf{databyte 0 i den mest signifikanta positionen}. Det är samma
packning på båda sidor av ledningen.

\section{Vad det kostar i tid}
\label{sec:timing}

En transaktion är 5 bytes = 40 bitar. Vid 1~MHz tar klockningen 40~µs, plus \code{SS}-hantering.
Räkna med \textbf{40--50~µs per transaktion}:

\begin{narrowtable}{@{}lrl@{}}
\thead{Operation} & \thead{Transaktioner} & \thead{Ungefär} \\ \midrule
\code{hasError()} & 1 & 45~µs \\
\code{receive()}, ingen frame & 1 & 45~µs \\
\code{receive()}, en frame & 6 & 270~µs \\
\code{send()} & 6 & 270~µs \\
\end{narrowtable}

Sätt det mot CAN-sidan: en maximal CAN-frame är omkring 130 bitar, alltså \textbf{130~µs} vid
1~Mbit/s.

\begin{limitation}
\textbf{En \code{send()} tar längre tid än framen den skickar.} SPI-länken, inte CAN-bussen, är
flaskhalsen i den här konstruktionen.

\textbf{Och pollningen kan missa frames.} Kontrollern buffrar inte: kommer två frames tätt inpå
varandra och din \code{receive()} tar 270~µs, hinner den andra skriva över den första. Det är
hårdvarans dokumenterade begränsning, inte din bugg --- men det är din kod som märker den.
\end{limitation}
